% Part: many-valued-logic
% Chapter: infinite-valued-logics
% Section: goedel

\documentclass[../../../include/open-logic-section]{subfiles}

\begin{document}

\olfileid{mvl}{inf}{god}

\olsection{ગોડેલનાં તર્કશાસ્ત્રો}

\begin{editorial}
  અનંતમૂલ્ય ગોડેલ તર્કશાસ્ત્ર અંગેનો આ વિભાગ ટૂંકું
  ``પ્રારૂપ'' માત્ર છે.
\end{editorial}

\begin{defn}\ollabel{def:goedel} અનંતમૂલ્ય ગોડેલ
તર્કશાસ્ત્ર~$\LogGod[\infty]$ નીચેની શ્રેણિક વડે વ્યાખ્યાયિત થાય છે:
\begin{enumerate}
  \item $\lfalse$, $\lnot$, $\land$, $\lor$, $\lif$ ધરાવતી
  પ્રમાણભૂત વિધાનાત્મક ભાષા~$\Lang L_0$.
  \item સત્યમૂલ્યોનો ગણ $V_\infty$.
  \item એકમાત્ર નિર્દિષ્ટ મૂલ્ય $1$ છે; એટલે $V^+ = \{1\}$.
  \item સત્યવિધેયો નીચેના વિધેયો વડે અપાય છે:
  \begin{align*}
    \tf{\lfalse} & = 0\\
    \tf{\lnot}[\LogGod](x) & = \begin{cases}
      1 & \text{જો } x =0\\
      0 & \text{અન્યથા}
    \end{cases}\\
    \tf{\land}[\LogGod](x,y) & = \min(x,y)\\
    \tf{\lor}[\LogGod](x,y) & = \max(x,y)\\
    \tf{\lif}[\LogGod](x,y) & = \begin{cases}
      1 & \text{જો } x \le y\\
      y & \text{અન્યથા.}
    \end{cases}
    \end{align*}
\end{enumerate}
$m$-મૂલ્ય ગોડેલ તર્કશાસ્ત્રની વ્યાખ્યા પણ આવી જ છે;
માત્ર $V = V_m$ છે.
\end{defn}
\sourcecorrection{OLMV-015}{સ્થિર અંગ્રેજી મૂળના
નિષેધ-વિધેયની \texttt{cases} રચના પહેલેથી ગણિતીય
વાતાવરણમાં હોવા છતાં પરિણામો આસપાસ ફરી
\texttt{\$} ચિહ્ન મૂકે છે. તે અયોગ્ય ગણિતીય
વાતાવરણ દૂર કરીને $1$ અને $0$નાં મૂલ્યો
યથાવત રાખ્યાં છે.}

\begin{prop}
  \olref[thr][god]{defn:goedel} વડે વ્યાખ્યાયિત
  $\LogGod[3]$ એ \olref{def:goedel} વડે વ્યાખ્યાયિત
  $\LogGod[3]$ જેવું જ છે.
\end{prop}

\begin{proof}
  \olref[thr][god]{defn:goedel}માં આપેલી સંયોજકોની
  સત્યતાસારણીઓની સરખામણી \olref{def:goedel}નાં
  સમીકરણોથી મળતી સત્યતાસારણીઓ સાથે કરીને આ જોઈ શકાય છે:
  \begin{center}
    \begin{tabular}{c|c} 
      $\tf{\lnot}[\LogGod[3]]$ & \\ 
      \hline  
      $1$ & $0$ \\ 
      $1/2$ & $0$ \\
      $0$ & $1$ 
    \end{tabular}
    \quad
    \begin{tabular}{c|ccc} 
      $\tf{\land}[\LogGod]$ & $1$ & $1/2$ & $0$ \\ 
      \hline 
      $1$ & $1$ & $1/2$ & $0$ \\ 
      $1/2$ & $1/2$ & $1/2$ & $0$\\ 
      $0$ & $0$ & $0$ & $0$ 
    \end{tabular}
    \\[2ex]
    \begin{tabular}{c|ccc} 
      $\tf{\lor}[\LogGod]$ & $1$ & $1/2$ & $0$ \\ 
      \hline 
      $1$ & $1$ & $1$ & $1$ \\ 
      $1/2$ & $1$ & $1/2$ & $1/2$ \\
      $0$ & $1$ & $1/2$ & $0$ 
    \end{tabular}
    \quad
    \begin{tabular}{c|ccc} 
      $\tf{\lif}[\LogGod]$ & $1$ & $1/2$ & $0$ \\ 
      \hline 
      $1$ & $1$ & $1/2$ & $0$ \\ 
      $1/2$ & $1$ & $1$ & $0$  \\ 
      $0$ & $1$ & $1$ & $1$ 
    \end{tabular}
  \end{center} 
\end{proof}

\begin{prop}\ollabel{prop:god-infty-m}
  જો $\Gamma \Entails[\LogGod[\infty]] !B$, તો બધાં~$m \ge 2$
  માટે $\Gamma \Entails[\LogGod[m]] !B$.
\end{prop}

\begin{proof}
  અભ્યાસપ્રશ્ન.
\end{proof}

\begin{prob}
  \olref[mvl][inf][god]{prop:god-infty-m} સાબિત કરો.
\end{prob}

હકીકતમાં, સાન્ત આધારગણ માટે વિપરીત દિશા પણ સાચી છે.
\sourcecorrection{OLMV-016}{સ્થિર અંગ્રેજી મૂળ
વિપરીત દિશા કોઈ મર્યાદા વિના કહે છે. સાન્ત
આધારગણ માટે સંમેય વિરોધી મૂલ્યાંકનમાં આવતા
મર્યાદિત સત્યમૂલ્યો એક જ સાન્ત જાળીમાં સમાય છે.
પરંતુ અનંત આધારગણ ગોડેલના પ્રેરણ વડે અસીમ
ઉતરતા મૂલ્યો ફરજિયાત કરી શકે; તે દરેક સાન્ત
જાળીમાં નિષ્કર્ષને નિર્દિષ્ટ કરે છતાં અનંત
સંમેય જાળીમાં વિરોધી મૂલ્યાંકન આપે છે.
તેથી અહીં દાવો સાન્ત આધારગણ સુધી મર્યાદિત કર્યો છે.}

$\LogGod[3]$ની જેમ $\LogGod[\infty]$માં બધી સ્ફુરણાવાદી રીતે
પ્રામાણ્ય !!{formula}s પુનરુક્તિઓ છે, અને નીચે દર્શાવેલી
બિનપુનરુક્તિઓ~$\LogGod[\infty]$માં પણ બિનપુનરુક્તિઓ છે:
\begin{align*}
  & p \lor \lnot p && (p \lif q) \lif (\lnot p \lor q) \\
  & \lnot\lnot p \lif p && \lnot(\lnot p \land \lnot q) \lif (p \lor q) \\
  & ((p \lif q) \lif p) \lif p && \lnot(p \lif q) \lif (p \land \lnot q)
\end{align*}
સ્ફુરણાવાદી રીતે અપ્રામાણ્ય !!{formula} હોવા છતાં $\LogGod[3]$ની પુનરુક્તિ
હોવાનું ઉદાહરણ $(p \lif q) \lor (q \lif p)$ એ
$\LogGod[\infty]$માં પણ પુનરુક્તિ છે. હકીકતમાં,
$\LogGod[\infty]$ને $(!A \lif !B) \lor (!B \lif !A)$ની
યોજના ઉમેરેલા સ્ફુરણાવાદી તર્કશાસ્ત્ર તરીકે લક્ષણિત
કરી શકાય છે. આ માઇકલ ડમેટે બતાવ્યું હતું, અને તેથી
$\LogGod[\infty]$ને ગોડેલ--ડમેટ તર્કશાસ્ત્ર~$\Log{LC}$ પણ કહે છે.

\begin{prob}
  $(p \lif q) \lor (q \lif p)$ એ~$\LogGod[\infty]$ની
  પુનરુક્તિ છે તે બતાવો.
\end{prob}

\begin{prob}
  $(p \lif q) \lor (q \lif r) \lor (r \lif s)$ એ
  $\LogGod[3]$ની પુનરુક્તિ છે, પરંતુ~$\LogGod[\infty]$ની
  પુનરુક્તિ નથી, તે બતાવો.
\end{prob}

\end{document}
